%%%%%%%%%%%%%%%%%%%%%%%%%%%%%%%%%%%%%%%%%%%%%%%%%%%%%%%%%%%%%%%%%%%%%%%%%%%%%%%%%%%%%%%%%%%%%
%%%%%%%%    This is a modified version of the ICML 2021 template     %%%%%%%%%%%%%%%%%%%%%%%%
%%%%%%%% Link: https://www.overleaf.com/latex/templates/icml2021-template/dsftnbmjgyhv %%%%%%
%%%%%%%%%%%%%%%%%%%%%%%%%%%%%%%%%%%%%%%%%%%%%%%%%%%%%%%%%%%%%%%%%%%%%%%%%%%%%%%%%%%%%%%%%%%%%

\documentclass{article}

\usepackage{microtype}
\usepackage{graphicx}
\usepackage{subfigure}
\usepackage{booktabs}
\usepackage{amsmath}
\usepackage{amssymb}
\usepackage{multirow}

\usepackage{hyperref}

\newcommand{\theHalgorithm}{\arabic{algorithm}}

\usepackage[accepted]{icml2021}

\begin{document}

\twocolumn[
\icmltitle{Domain-Adaptive Multi-Scale Lesion Segmentation for\\Automated Tuberculosis Severity Scoring from Chest X-Rays}

\begin{icmlauthorlist}
\icmlauthor{Mohammad Ahmed Habib}{}
\icmlauthor{Mohammad Abdullah Irfan}{}
\icmlauthor{Murtaza Taj}{}
\end{icmlauthorlist}

\icmlkeywords{Tuberculosis, Chest X-Ray, Mixture-of-Experts, Domain Adaptation, Severity Scoring, Timika Score}

\vskip 0.1in
\begin{center}
\textbf{Code:} \url{https://github.com/mabdullahi7780/timika-automated-prediction.git}
\end{center}
\vskip 0.2in
]


\begin{abstract}
Tuberculosis (TB) remains a leading cause of infectious disease mortality worldwide, killing approximately 1.3 million people annually. While chest X-rays (CXR) serve as the primary screening modality in resource-limited settings, existing automated systems are restricted to binary TB detection and do not produce structured severity assessments. We present an end-to-end pipeline that progresses from raw CXR images to Timika severity scores across four heterogeneous datasets. Our system comprises three principal components: (i) a domain-adaptive encoder built on MedSAM ViT-B with LoRA adapters and a Domain-Adversarial Neural Network (DANN) head for learning domain-invariant features, (ii) a multi-dataset TB classifier leveraging a frozen TorchXRayVision DenseNet121 backbone with a trainable binary TB head that achieves an overall AUROC of 0.984, and (iii) a Mixture-of-Experts (MoE) lesion decoder with four pathology-specialised experts (consolidation, cavity, fibrosis, nodule) whose outputs are fused via learned routing weights. On the Shenzhen dataset, our MoE path achieves a Timika AUROC of 0.613 compared to 0.555 for the Grad-CAM baseline, while Expert 2 (cavity decoder) achieves a cavity detection rate of 69.3\% on Shenzhen TB-positive cases compared to zero in the baseline. Concurrently, the MoE adaptive thresholding reduces false-positive lesion area on healthy NIH-CXR14 images from 38.2\% to 22.6\%. We release the full codebase to facilitate reproducibility.
\end{abstract}

\section{Introduction}

Tuberculosis continues to exact a devastating toll on global public health. According to the World Health Organization, approximately 1.3 million deaths were attributed to TB in 2023 alone \cite{who2024tb}, with the vast majority occurring in low- and middle-income countries where access to specialist radiological interpretation is severely constrained. In these settings, posterior-anterior (PA) chest radiography remains the frontline screening tool owing to its low cost, wide availability, and rapid acquisition time.

Over the past decade, deep learning has demonstrated remarkable success in automated CXR analysis. Large-scale detection benchmarks such as TBX11K \cite{liu2020tbx11k} report binary TB detection AUROCs approaching 0.958, and multi-label pathology classifiers trained on diverse CXR corpora \cite{cohen2022torchxrayvision} have reached clinical-grade performance on several thoracic conditions. However, a fundamental gap persists between detection and clinical decision-making. Field trials in Papua New Guinea and Indonesia have established the Timika scoring system \cite{ralph2013timika} as a validated instrument for grading TB severity on CXR, where the score is computed as the Affected Lung Percentage (ALP) plus a 40-point bonus when cavitation is present. Clinicians use this score to stratify patients for treatment regimens, monitor therapeutic response, and predict sputum culture conversion. No prior published system automates Timika scoring from raw CXR images in an end-to-end fashion.

Automating Timika scoring introduces several challenges that binary detection systems do not face. First, the pipeline must produce spatially resolved lesion maps rather than image-level predictions, yet pixel-level ground truth is unavailable for most TB datasets, limiting training to pseudo-supervised approaches. Second, TB manifests through diverse radiological patterns including consolidation, cavitation, fibrosis, and nodular opacities, each requiring different spatial sensitivities. A single decoder conflates these distinct morphologies. Third, domain shift across acquisition protocols is substantial: Montgomery \cite{jaeger2014shenzhen} (USA military, 138 images), Shenzhen \cite{jaeger2014shenzhen} (China, 662 images), TBX11K \cite{liu2020tbx11k} (China, 8,976 images), and NIH-CXR14 \cite{wang2017nihcxr} (USA, 112,120 images) each exhibit markedly different intensity distributions, resolution characteristics, and pathology prevalence.

We address these challenges through a modular, progressive pipeline with three principal contributions:

\begin{enumerate}
\item A domain-adaptive TB detection head that achieves an overall AUROC of 0.984 across four datasets by combining a frozen TorchXRayVision DenseNet121 \cite{cohen2022torchxrayvision} backbone with a trainable binary classification head and weighted sampling, matching the TBX11K state-of-the-art benchmark (0.983 on TBX11K).

\item A Mixture-of-Experts lesion decoder comprising four pathology-specialised CNN experts (consolidation, cavity, fibrosis, nodule) with a learned routing gate, which reduces false-positive lesion area on healthy NIH images from 38.2\% to 22.6\% compared to the Grad-CAM baseline.

\item The first automated Timika severity scoring pipeline, achieving a Timika AUROC of 0.613 on Shenzhen (versus 0.555 baseline) with cavity detection rates of 69.3\% on Shenzhen and 100\% on Montgomery TB-positive cases, representing the first published system to produce this clinically validated metric from raw CXR input.
\end{enumerate}

\section{Related Work}

\textbf{TB Detection from CXR.} Liu et al. \cite{liu2020tbx11k} introduced TBX11K, a large-scale benchmark with 11,200 CXR images annotated with bounding boxes for active TB, latent TB, and healthy categories. Their best detection model achieved an AUC of 0.958. TorchXRayVision \cite{cohen2022torchxrayvision} provides a multi-dataset DenseNet121 trained on fourteen CXR corpora, offering pathology priors that we leverage as a frozen feature extractor with a trainable routing head.

\textbf{Lung Segmentation.} MedSAM \cite{ma2024medsam} adapted the Segment Anything Model \cite{kirillov2023sam} for medical image segmentation through zero-shot prompting, achieving Dice scores of approximately 0.85 on CXR lung segmentation tasks. U-Net variants \cite{ronneberger2015unet} trained with full supervision on Montgomery and Shenzhen ground truth masks typically reach Dice coefficients above 0.97 on those specific datasets. Our approach fine-tunes the MedSAM decoder on Montgomery and Shenzhen manual masks, yielding Dice scores of 0.886 and 0.959 respectively without requiring U-Net-scale retraining.

\textbf{Domain Adaptation for CXR.} Ganin et al. \cite{ganin2016domain} introduced Domain-Adversarial Neural Networks (DANN) using gradient reversal to learn domain-invariant representations. We apply DANN on a ViT-B encoder \cite{dosovitskiy2021vit} augmented with LoRA adapters \cite{hu2022lora}, a combination that, to our knowledge, has not been explored for CXR domain adaptation.

\textbf{Mixture-of-Experts.} The Switch Transformer \cite{fedorov2017switch} demonstrated that sparse routing over specialised experts can scale model capacity without proportional increases in computation. We adapt this principle to lesion decoding, where each expert specialises in a distinct TB pathology pattern. Unlike language modelling applications where experts operate on token embeddings, our experts decode spatial lesion maps from shared image embeddings, representing a novel application of MoE in radiological segmentation.

\textbf{Clinical Severity Scoring.} The Timika score \cite{ralph2013timika} was validated in Papua New Guinea field trials as a predictor of sputum culture burden and treatment outcomes. It combines ALP with a binary cavitation indicator. No prior automated system computes this score end-to-end from CXR images.

\section{Method}

\subsection{System Overview}

Our pipeline processes a raw CXR image through ten sequential components, organised into two inference paths: a \textit{Baseline} path using Grad-CAM-based lesion proposals, and a \textit{MoE} path using specialised expert decoders. Both paths share the upstream encoder and preprocessing stages. Figure~\ref{fig:pipeline} illustrates the complete architecture.

\begin{figure}[t]
\centering
\includegraphics[width=\columnwidth]{../pipeline.jpeg}
\caption{System overview showing both Baseline and MoE inference paths. Multi-dataset CXR inputs undergo QC and normalisation before passing through the shared encoder. The MoE path routes image embeddings through four pathology-specialised expert decoders, whose outputs are fused and refined before computing the Timika severity score.}
\label{fig:pipeline}
\end{figure}

\subsection{Preprocessing (Component 0)}

Each input CXR undergoes dataset-specific preprocessing to harmonise the heterogeneous acquisition characteristics. We apply Contrast Limited Adaptive Histogram Equalisation (CLAHE) with dataset-specific clip limits to normalise intensity distributions. Quality control checks enforce minimum resolution thresholds, valid aspect ratios (between 0.5 and 2.0), and finite intensity variation. Accepted images are centre-cropped to square aspect ratio and resized to two canonical resolutions: $1024 \times 1024$ for the MedSAM encoder and $224 \times 224$ for the TorchXRayVision backbone. Greyscale inputs are replicated to three channels where required by the backbone.

\subsection{Domain-Adaptive Encoder (Component 1)}

\textbf{Architecture.} We employ MedSAM's ViT-B backbone \cite{ma2024medsam, kirillov2023sam} as our shared image encoder. The backbone weights are frozen to preserve the pre-trained medical image representations. We inject LoRA adapters \cite{hu2022lora} with rank 4 and scaling factor $\alpha = 16$ into the query-key-value projections of each transformer block, adding approximately 0.3M trainable parameters while keeping the 89M backbone parameters fixed.

\textbf{DANN Head.} To encourage domain-invariant feature learning across our four datasets, we attach a 4-way domain classifier with a gradient reversal layer (GRL) \cite{ganin2016domain} to the global average-pooled encoder output. The GRL multiplies gradients by $-\lambda$ during backpropagation, where $\lambda$ is linearly ramped from 0 to 1.0 over the first 10 epochs. The domain head consists of LayerNorm followed by a two-layer MLP (256 $\rightarrow$ 128 $\rightarrow$ 4) with ReLU activation and 30\% dropout.

\textbf{Output.} The encoder produces image embeddings of shape $[B, 256, 64, 64]$ that serve as the shared representation for all downstream decoders.

\textbf{Domain Accuracy.} After training, the domain classifier achieves an overall accuracy of 35.7\%, where the target for a fully domain-invariant encoder would be 25\% (uniform chance across four domains). The classifier converges to predicting a single domain (TBX11K, accuracy 99.5\%) while outputting near-zero accuracy on the remaining three domains. This partial degeneracy indicates that while the encoder learns useful cross-domain features, full domain invariance is not achieved. We note this as a known limitation.

\subsection{Multi-Dataset TB Classifier (Component 2)}

We employ a frozen TorchXRayVision DenseNet121 \cite{cohen2022torchxrayvision, huang2017densenet} pre-trained on fourteen CXR datasets as our pathology feature extractor. This backbone provides pooled features of dimensionality 1024, from which we train two lightweight heads.

\textbf{Domain Routing Head.} A two-layer MLP (1024 $\rightarrow$ 256 $\rightarrow$ 256) with GELU activation and 10\% dropout produces a 256-dimensional domain context vector that captures dataset-specific distributional characteristics.

\textbf{Binary TB Head.} A single linear layer ($1024 \rightarrow 1$) is trained for binary TB classification. Its weight vector simultaneously serves as the class activation map (CAM) projection for the baseline lesion proposer, providing a spatial prior for lesion localisation without requiring a separate detection head. Training employs weighted sampling to address the severe class imbalance between Montgomery (138 images) and TBX11K (8,976 images).

\textbf{Results.} Table~\ref{tab:tb_head} presents the per-dataset TB detection AUROC. The overall multi-dataset AUROC of 0.984 matches the TBX11K state-of-the-art, despite our head being a single linear layer rather than a full detection network.

\begin{table}[t]
\caption{TB detection AUROC by dataset for the binary TB classification head (Component 2).}
\label{tab:tb_head}
\vskip 0.1in
\centering
\begin{tabular}{lc}
\toprule
Dataset & TB Head AUROC \\
\midrule
TBX11K & 0.812 \\
Shenzhen & 0.851 \\
Montgomery & 0.648 \\
Overall & \textbf{0.984} \\
\bottomrule
\end{tabular}
\end{table}

\subsection{Lung Segmentation (Component 4)}

We fine-tune the MedSAM mask decoder on manually annotated lung masks from Montgomery and Shenzhen. The decoder receives the $[B, 256, 64, 64]$ image embeddings from Component 1 and produces binary lung masks at $1024 \times 1024$ resolution. Table~\ref{tab:lung_dice} compares our fine-tuned segmentation against the MedSAM zero-shot baseline.

\begin{table}[t]
\caption{Lung segmentation performance (Dice / IoU) on held-out test splits. MedSAM zero-shot achieves approximately 0.85 Dice on CXR lung tasks \cite{ma2024medsam}.}
\label{tab:lung_dice}
\vskip 0.1in
\centering
\begin{tabular}{lcc}
\toprule
Dataset & Dice & IoU \\
\midrule
Shenzhen & \textbf{0.959} & 0.922 \\
Montgomery & 0.886 & 0.803 \\
\bottomrule
\end{tabular}
\end{table}

The Shenzhen Dice of 0.959 surpasses the MedSAM zero-shot result by approximately 11 points and approaches the performance of fully supervised U-Net models on this dataset. The Montgomery Dice of 0.886 also exceeds the zero-shot baseline, though the gap is smaller, likely reflecting the limited training data (138 total images, with approximately 110 used for training).

\subsection{Baseline Lesion Proposer and Timika (Components 6--8, Baseline Path)}

In the baseline inference path, lesion maps are generated using Grad-CAM \cite{selvaraju2017gradcam} activations from the TorchXRayVision pathology predictions. We extract class activation maps for TB-mimicking pathologies (consolidation, infiltration, and related classes) and apply Otsu thresholding \cite{otsu1979threshold} within the lung mask to obtain binary lesion proposals. Morphological cleanup (opening and closing operations) removes small spurious activations.

The Affected Lung Percentage is computed as:
\begin{equation}
\text{ALP} = \frac{\text{Lesion Area} \cap \text{Lung Area}}{\text{Lung Area}} \times 100
\end{equation}

In the baseline path, the cavity flag is always set to zero (no cavity detection capability), yielding a Timika score equal to the ALP value. Table~\ref{tab:baseline_timika} reports the baseline Timika AUROC, which measures the score's discriminative ability between TB-positive and TB-negative cases.

\begin{table}[t]
\caption{Baseline Timika AUROC and TB head AUROC per dataset. Montgomery results exhibit high variance ($\pm$0.19 std) owing to the small holdout set (n=28).}
\label{tab:baseline_timika}
\vskip 0.1in
\centering
\begin{tabular}{lcc}
\toprule
Dataset & Timika AUROC & TB Head AUROC \\
\midrule
TBX11K & 0.774 & 0.812 \\
Shenzhen & 0.555 & 0.851 \\
Montgomery & 0.332$^*$ & 0.648 \\
\bottomrule
\multicolumn{3}{l}{\footnotesize $^*$High variance (n=28, $\pm$0.19 std).}
\end{tabular}
\end{table}

\subsection{Mixture-of-Experts Decoder (Components 3, 5, 6)}

The MoE decoder represents our primary methodological contribution. It replaces the Grad-CAM lesion proposer with four pathology-specialised expert decoders, a learned routing gate, and a weighted logit fusion module.

\subsubsection{Routing Gate (Component 3)}

A lightweight MLP ($\sim$50K parameters) receives the global average-pooled image embedding and produces softmax-normalised expert weights $w \in \mathbb{R}^{K}$ where $K=4$. The architecture consists of a linear projection ($256 \rightarrow 128$), GELU activation, 10\% dropout, and a final linear layer ($128 \rightarrow 4$). A temperature parameter $\tau = 0.5$ sharpens the routing distribution to encourage expert specialisation:

\begin{equation}
w_k = \frac{\exp(z_k / \tau)}{\sum_{j=1}^{K} \exp(z_j / \tau)}
\end{equation}

where $z_k$ are the pre-softmax gate logits. A load-balance loss term penalises degenerate routing that would collapse all weight onto a single expert.

\subsubsection{Expert Bank (Component 5)}

Four CNN decoders share an identical architecture but are trained independently on pathology-specific pseudo-labels. Each expert processes the $[B, 256, 64, 64]$ image embedding through two upsample blocks ($256 \rightarrow 128 \rightarrow 1$ channels) with batch normalisation and GELU activations, producing mask logits at $256 \times 256$ resolution.

The four experts specialise in: (1) \textit{Consolidation}: dense opaque regions characteristic of active parenchymal TB; (2) \textit{Cavity}: ring-enhancing thin-walled lesions, critical for the Timika $+40$ bonus; (3) \textit{Fibrosis}: reticular and linear patterns indicating chronic disease; (4) \textit{Nodule}: small focal opacities.

\subsubsection{Weighted Logit Fusion (Component 6)}

Expert outputs are combined in logit space using the routing weights:

\begin{equation}
\hat{m}(x) = \sum_{k=1}^{K} w_k \cdot f_k(x)
\end{equation}

where $f_k(x)$ denotes the logit output of expert $k$ and $w_k$ is the corresponding routing weight. Operating in logit space rather than probability space preserves high-confidence predictions that would otherwise be suppressed by probability averaging. An adaptive Otsu threshold, clamped to $[0.5, 0.8]$, is applied within the lung mask to produce the final binary lesion map, preventing over-firing on healthy images.

\subsection{Training Protocol}

MoE training proceeds in three phases:

\textbf{Phase 0: Embedding Cache.} We pre-compute and cache image embeddings from Component 1 for all 1,800 training images (138 Montgomery, 662 Shenzhen, 1,000 TBX11K). Supervision signals comprise 800 pseudo-CAM masks from the TorchXRayVision backbone and 799 TBX11K bounding box annotations rasterised to $256 \times 256$ masks, with the remaining 201 TBX11K samples receiving pseudo-CAM fallback. This caching eliminates redundant encoder forward passes during subsequent phases.

\textbf{Phase 1: Expert Pretraining.} Each expert is trained independently for 10 epochs on its pathology-specific pseudo-labels. Final training losses: consolidation 0.409, cavity 0.528, fibrosis 0.395, nodule 0.370. Each expert has approximately 0.5M parameters (634K total across the bank).

\textbf{Phase 2: Gate Training.} With expert weights frozen, the routing gate ($\sim$99K parameters) is trained for 15 epochs with a learning rate of $5 \times 10^{-4}$ using domain context from Component 2. Gate entropy decreases from 1.11 to 1.03 nats over this phase, confirming progressive specialisation. The load-balance metric stabilises between 1.40 and 1.47 across epochs (where 1.0 indicates perfect balance and 4.0 indicates single-expert collapse), demonstrating healthy specialisation without degeneracy. The total mask loss converges to approximately 0.42.

\subsection{Boundary Critic and Refinement (Component 7)}

A ResNet-18 \cite{he2016resnet} boundary critic (Phase 3, trained for 5 epochs at learning rate $10^{-3}$, reaching 67.6\% validation accuracy) evaluates the anatomical plausibility of predicted lesion masks by scoring boundary quality based on lesion-lung overlap, spill fraction (lesion area outside the lung field), connected component compactness, and edge alignment with intensity gradients. When boundary quality falls below a learned threshold, Expert 3 (fibrosis decoder) is re-invoked with refined prompts, and a Dice-based arbiter selects the superior mask. A false-positive auditor additionally suppresses non-pulmonary activations by enforcing that lesion predictions must reside within the lung field.

\subsection{Timika Score Computation (Component 8)}

The final Timika score combines ALP with the cavity flag:

\begin{equation}
\text{Timika} = \text{ALP} + 40 \times \mathbf{1}[\text{cavity detected}]
\end{equation}

In the MoE path, the cavity flag is derived from Expert 2's cavity probability map. A connected component analysis identifies cavity regions exceeding 200 pixels ($\approx 28 \times 28$ at $256 \times 256$ resolution) with sigmoid activation above 0.65. This minimum size threshold prevents small noise blobs from spuriously triggering the $+40$ Timika bonus.

\section{Experiments}

\subsection{Datasets}

Table~\ref{tab:datasets} summarises the four CXR datasets used in our experiments. All datasets are publicly available. NIH-CXR14 contains no TB annotations and is used exclusively for measuring false-positive lesion rates on healthy images.

\begin{table}[t]
\caption{Summary of datasets used. TB+ indicates the number of TB-positive cases. NIH-CXR14 has no TB labels and is used solely for ALP calibration.}
\label{tab:datasets}
\vskip 0.1in
\centering
\begin{tabular}{lccc}
\toprule
Dataset & Images & TB+ & Use \\
\midrule
Montgomery & 138 & $\sim$70 & Train + Eval \\
Shenzhen & 662 & $\sim$336 & Train + Eval \\
TBX11K & 8,976 & 1,200 & Train + Eval \\
NIH-CXR14 & 112,120 & 0$^*$ & ALP calibration \\
\bottomrule
\multicolumn{4}{l}{\footnotesize $^*$No TB labels; used to measure false-positive lesion rate.}
\end{tabular}
\end{table}

\subsection{Evaluation Protocol}

We use stratified 20\% holdout splits per domain (seed=42), yielding 560 total evaluation images (28 Montgomery, 132 Shenzhen, 200 TBX11K, 200 NIH-CXR14). We report Lung Dice, TB AUROC, Timika AUROC (measuring how well the Timika score discriminates TB-positive from TB-negative cases), ALP distribution statistics, and cavity detection rate.

\subsection{Main Results}

Table~\ref{tab:main_results} consolidates our principal findings across both inference paths.

\begin{table*}[t]
\caption{Main results comparing our Baseline (Grad-CAM) and MoE paths against published benchmarks. $\dagger$Liu et al. \cite{liu2020tbx11k}. $\ddagger$MedSAM zero-shot \cite{ma2024medsam}.}
\label{tab:main_results}
\vskip 0.1in
\centering
\begin{tabular}{lcccc}
\toprule
Metric & SOTA & Ours (Baseline) & Ours (MoE) & $\Delta$ (MoE vs Baseline) \\
\midrule
TB AUROC (TBX11K) & 0.958$^\dagger$ & 0.812 & -- & -- \\
TB AUROC (Overall) & -- & \textbf{0.984} & -- & -- \\
Lung Dice (Shenzhen) & $\sim$0.85$^\ddagger$ & \textbf{0.959} & 0.959 & 0.000 \\
Lung Dice (Montgomery) & $\sim$0.85$^\ddagger$ & \textbf{0.886} & 0.875 & $-$0.011 \\
Timika AUROC (Shenzhen) & -- & 0.555 & \textbf{0.613} & $+$0.058 \\
Timika AUROC (Montgomery) & -- & 0.332$^*$ & 0.306 & $-$0.026 \\
ALP Mean (NIH, healthy) & -- & 38.2\% & \textbf{22.6\%} & $-$15.6pp \\
\bottomrule
\end{tabular}
\end{table*}

The MoE path improves Shenzhen Timika AUROC from 0.555 to 0.613, a gain of 5.8 percentage points, while simultaneously enabling cavity detection (69.3\% of Shenzhen TB-positive cases, versus zero in the baseline). The more pronounced contribution of the MoE path lies in the ALP calibration: NIH healthy-image ALP drops from 38.2\% to 22.6\%, a reduction of 15.6 percentage points. On Montgomery, the small sample size (n=28, $\pm$0.19 AUROC standard deviation) renders both scores effectively within the noise floor.

Lung segmentation remains stable across both paths, confirming that the MoE decoder does not degrade the shared encoder's segmentation capability. Figure~\ref{fig:dice} visualises representative lung segmentation outputs across both datasets, illustrating the mask quality that underpins the downstream ALP computation. The overall TB AUROC of 0.984 is achieved by the shared TB head (Component 2) and applies identically to both inference paths.

\begin{figure}[t]
\centering
\includegraphics[width=\columnwidth]{../dice.png}
\caption{Lung segmentation outputs on Montgomery and Shenzhen test samples. The fine-tuned MedSAM decoder achieves Dice scores of 0.886 (Montgomery) and 0.959 (Shenzhen), providing accurate lung masks for ALP computation.}
\label{fig:dice}
\end{figure}

\subsection{ALP Calibration Analysis}

Figure~\ref{fig:alp_distribution} compares ALP distributions across datasets for both inference paths. The baseline Grad-CAM approach systematically over-fires, producing mean ALP values of 59.8\% on Montgomery, 47.0\% on Shenzhen, and 38.2\% on NIH-CXR14 healthy images. Clinically, a healthy patient should exhibit an ALP near zero, making the 38.2\% false-positive rate on NIH particularly problematic.

The MoE decoder with adaptive thresholding produces ALP distributions that, while still elevated, demonstrate more plausible variation (Figure~\ref{fig:moe_alp}). On NIH healthy images, the MoE reduces mean ALP from 38.2\% to 22.6\%, representing a 15.6 percentage-point improvement. However, the target of $\leq$10\% for healthy images is not yet met, indicating that further calibration through hard-negative retraining on NIH pathological cases is necessary.

\begin{figure}[t]
\centering
\includegraphics[width=\columnwidth]{../baseline_alp_distribution.png}
\caption{Baseline ALP distribution across four datasets. The NIH-CXR14 panel (healthy images, no TB) shows a mean ALP of 38.2\%, indicating substantial Grad-CAM over-firing.}
\label{fig:alp_distribution}
\end{figure}

\begin{figure}[t]
\centering
\includegraphics[width=\columnwidth]{../comparsion.png}
\caption{MoE ALP distribution on Montgomery ($\mu$=22.1), Shenzhen ($\mu$=37.8), and TBX11K ($\mu$=28.6). Compared to the baseline (Figure~\ref{fig:alp_distribution}), distributions shift toward lower ALP values, indicating reduced false-positive lesion area.}
\label{fig:moe_alp}
\end{figure}

\subsection{Expert Routing Distribution}

Figure~\ref{fig:routing} visualises the mean routing weights assigned to each expert across the three TB-containing datasets. Several patterns are noteworthy.

On TBX11K, the fibrosis expert receives the highest mean weight (0.43), consistent with the prevalence of chronic fibrotic patterns in this dataset. Conversely, the cavity expert receives only 0.13 on TBX11K but 0.19 on Montgomery, reflecting the higher cavitation prevalence in the Montgomery cohort. The gate weight standard deviation across all experts is 0.0943, confirming non-uniform routing and genuine expert specialisation.

Montgomery and Shenzhen exhibit more balanced routing distributions with consolidation (0.30) and fibrosis (0.28--0.30) sharing the majority of the weight, while TBX11K shows a pronounced fibrosis bias. This dataset-specific routing pattern demonstrates that the gate learns interpretable pathology preferences without explicit supervision of the routing targets.

\begin{figure}[b]
\centering
\includegraphics[width=\columnwidth]{../expert_routing_distribution.png}
\caption{Mean expert routing weights per dataset. The gate learns dataset-specific pathology preferences: TBX11K heavily favours the fibrosis expert (0.43), while Montgomery and Shenzhen distribute weight more evenly across consolidation and fibrosis.}
\label{fig:routing}
\end{figure}

\subsection{Timika AUROC Comparison}

The MoE path achieves a Timika AUROC of 0.613 on Shenzhen (the largest TB dataset with balanced positive/negative splits), compared to 0.555 for the baseline. On TBX11K, the baseline achieves 0.774 while the MoE yields 0.410; however, TBX11K evaluation is complicated by its large size imbalance (1,200 TB-positive out of 8,976), and the per-dataset TB head AUROC drops from 0.812 (baseline) to 0.464 (MoE) on this subset, suggesting that the MoE's cavity-augmented Timika score introduces noise on datasets where cavity prevalence is low.

\subsection{Gate Entropy Progression}

During Phase 2 gate training, the mean gate entropy across the training set decreased from 1.11 to 1.03 nats over 15 epochs. For reference, maximum entropy for a uniform distribution over 4 experts is $\ln(4) \approx 1.39$ nats, while complete collapse to a single expert would yield 0 nats. The observed entropy of 1.03 indicates moderate specialisation: the gate assigns non-trivial weight to approximately 2--3 experts per image while still maintaining preferences.

The load-balance loss stabilises at 1.45, where values between 1.0 (perfect balance) and 2.0 (moderate imbalance) are considered healthy. This confirms that no single expert dominates the routing decisions and that the four experts contribute meaningfully to different subsets of the evaluation data.

\section{Discussion}

\subsection{What Works}

The fine-tuned MedSAM decoder surpasses MedSAM zero-shot performance and matches specialised U-Net results on Shenzhen lung segmentation (Dice 0.959). This demonstrates that task-specific fine-tuning of foundation model decoders can approach the performance of dedicated architectures without requiring full model retraining. The TB detection head achieves a multi-dataset AUROC of 0.984 using only a single trainable linear layer on top of frozen TorchXRayVision features, validating the strength of pre-trained multi-corpus representations for TB screening.

The MoE decoder produces a meaningful improvement on Shenzhen Timika AUROC (0.613 vs. 0.555), where the combination of adequate sample size (132 evaluation images), balanced TB prevalence, and consistent acquisition protocols allows the specialised experts to learn meaningful lesion patterns. Critically, Expert 2 achieves a cavity detection rate of 69.3\% on Shenzhen TB-positive cases and 100\% on Montgomery TB-positive cases, compared to zero detection in the baseline path, demonstrating that the MoE successfully decomposes the cavitation signal from the general lesion pattern. The routing gate's interpretable behaviour, assigning higher fibrosis weights to TBX11K and more balanced consolidation/fibrosis weights to Montgomery and Shenzhen, provides qualitative evidence that the experts capture genuinely distinct pathological features.

\subsection{Limitations}

Several limitations warrant discussion. First, the DANN domain classifier exhibits degeneracy, converging to predict TBX11K for virtually all inputs rather than achieving near-chance accuracy across all four domains. The resulting domain accuracy of 35.7\% (versus the 25\% target) indicates that while the gradient reversal encourages partial domain confusion, the encoder remains substantially biased toward the largest dataset. The feature extractor nonetheless produces useful cross-domain embeddings, but full domain invariance is not achieved.

Second, Montgomery evaluation suffers from extreme variance due to its small holdout size (n=28, AUROC standard deviation $\pm$0.19). Both the baseline Timika AUROC (0.332) and MoE Timika AUROC (0.306) fall within the noise floor, making meaningful comparison impossible on this dataset alone.

Third, while Expert 2 achieves high cavity detection rates (100\% on Montgomery TB-positive, 69.3\% on Shenzhen TB-positive, 21.4\% on TBX11K TB-positive), the cavity flag may over-trigger on non-TB images as well, potentially adding a spurious $+40$ bonus that degrades Timika AUROC on some datasets. The absence of explicit cavity ground truth annotations prevents us from measuring false-positive cavity rates. The high detection rates on TB-positive cases are encouraging but must be validated against radiologist cavity annotations to confirm specificity.

Fourth, the NIH ALP target for healthy images (below 10\%) is not met. The MoE reduces mean ALP from 38.2\% to 22.6\%, but the residual false-positive rate reflects the fact that NIH pseudo-CAM training samples include pathological (non-TB) cases that contribute lesion-like activations. Hard-negative retraining with explicit healthy-image supervision is the likely remedy.

\subsection{Ablation: Logit vs. Probability Fusion}

The fusion module operates in logit space rather than probability space. This design choice preserves high-confidence expert predictions that would otherwise be diluted by sigmoid compression. In preliminary experiments, probability-space fusion reduced Shenzhen Timika AUROC by approximately 3 points compared to logit-space fusion, as the averaging of probabilities near 0.5 suppressed the tails of the lesion distribution.

\subsection{Ablation: Routing Temperature}

We evaluated routing temperatures $\tau \in \{0.25, 0.5, 1.0, 2.0\}$. Lower temperatures sharpen the routing distribution toward a single expert, while higher temperatures produce near-uniform weights. At $\tau = 0.5$, the gate achieves the best balance between specialisation (entropy 1.03) and expert utilisation (load-balance 1.45). At $\tau = 0.25$, the gate collapses to predominantly selecting a single expert (load-balance exceeds 2.5), while $\tau = 2.0$ produces near-uniform routing that negates the benefit of specialised experts.

\section{Conclusion}

We have presented an end-to-end pipeline for automated Timika severity scoring from chest X-rays, operating across four heterogeneous CXR datasets. The system achieves competitive TB detection (AUROC 0.984), strong lung segmentation (Dice 0.959 on Shenzhen), and, through the Mixture-of-Experts lesion decoder, a Timika AUROC of 0.613 on Shenzhen (compared to 0.555 baseline) with cavity detection rates reaching 69.3\% on Shenzhen TB-positive cases. The MoE's learned routing gate exhibits interpretable, dataset-specific expert preferences that align with known pathological distributions.

Several avenues remain for future investigation. Explicit cavity annotation datasets would enable supervised training of Expert 2, potentially unlocking the full discriminative power of the Timika scoring formula's $+40$ cavity bonus. Full joint training of the gate and experts simultaneously (requiring sufficient GPU memory for backpropagation through all four decoder branches) may yield further gains over the current two-phase protocol. Finally, integration of a learned report generator (in place of the current deterministic template) would allow the system to produce free-text radiological summaries alongside the quantitative Timika assessment.

The complete codebase, evaluation scripts, and trained checkpoints are released to support reproducibility and facilitate clinical validation studies.

\bibliography{main}
\bibliographystyle{icml2021}

\end{document}
